% =====================================================================
% chapters/01_introduction.tex
% Chapter 1 — Introduction
% Project: Vision-Based Running Technique Classification from In-the-Wild
%          Internet Videos using Markerless Pose Estimation and Deep Learning
%
% No empirical metrics appear here. References use citation keys from the
% thesis bibliography.
% =====================================================================

\chapter{Introduction}
\label{ch:introduction}

% ---------------------------------------------------------------------
\section{Background}
\label{sec:intro-background}

Running technique --- the manner in which a person coordinates posture, foot
strike, stride, and limb motion through the running cycle --- has a direct
bearing on athletic performance, training efficiency, and the risk of
overuse injury \cite{novacheck1998biomechanics}. Coaches, clinicians, and runners
therefore have a long-standing interest in assessing whether a given running
form is sound or in need of correction. Distinguishing "correct" from "false"
technique is, however, a subtle judgement: the diagnostic information is
distributed across the whole kinematic chain and is expressed only as the body
moves periodically through repeated running cycles, rather than in any single
static posture \cite{inostroza2025runningcycle, cutting1978gaitperiodicity}.

Traditional approaches to running analysis are accurate but costly to apply at
scale. Optical marker-based motion capture provides laboratory-grade kinematics
but requires controlled environments, calibrated multi-camera rigs, and marker
placement; wearable inertial sensors enable field measurement but require
body-worn instrumentation and calibration; and expert visual inspection depends
on scarce specialist time and is inherently subjective
\cite{panconi2024markerless, bulling2014tutorial}. None of these
modalities can be applied retrospectively to footage that was not captured with
them in mind.

At the same time, \emph{in-the-wild} running videos --- clips recorded with
ordinary cameras and shared on the internet --- are abundant. This footage is an
enormous, untapped source of running data, but it is also noisy and
uncontrolled: it varies in camera viewpoint, resolution, frame rate, lighting,
and framing, and it carries no instrumentation or ground-truth measurement
\cite{topham2023hpesurvey}. Recent advances in \emph{markerless pose estimation}
make it possible to recover body landmarks directly from such ordinary RGB
video, with per-landmark confidence values, and thereby open the possibility of
analysing running technique from video that was never intended for biomechanical
study \cite{topham2023hpesurvey, lugaresi2019mediapipe, bazarevsky2020blazepose}. This thesis
investigates whether, and under what conditions, that possibility can be realised
with deep learning.

% ---------------------------------------------------------------------
\section{Problem Statement}
\label{sec:intro-problem}

The problem addressed in this thesis is stated as follows: \emph{given an
in-the-wild running video, classify the running technique as correct or false
using markerless pose estimation and deep learning.} The input is an ordinary
single-runner video with no instrumentation or calibration; the desired output
is a binary technique label, ideally accompanied by an indication of confidence.
The task must be solved without the controlled conditions assumed by
marker-based or sensor-based methods, relying instead on landmarks recovered
from the raw footage.

Solving this problem in the wild raises a set of inter-related challenges:

\begin{itemize}
  \item \textbf{Camera viewpoint variation.} Clips are filmed from arbitrary and
        inconsistent angles, so the same technique appears differently across
        videos.
  \item \textbf{Uncontrolled resolution and frame rate.} Spatial resolution and
        temporal sampling (FPS) differ between clips, affecting both pose
        quality and the temporal granularity of the running cycle.
  \item \textbf{Pose-landmark noise.} Markerless estimates from uncontrolled
        footage are noisier and less certain than laboratory measurements.
  \item \textbf{Missing or unstable lower-body landmarks.} The ankle, heel, and
        toe landmarks most diagnostic of running form are also the most
        frequently occluded, blurred, or mislocalised.
  \item \textbf{Foot-strike and cycle-boundary ambiguity.} The events that
        define a running cycle are themselves hard to localise reliably from
        noisy signals.
  \item \textbf{Automatic cycle detection affects classification.} Because the
        classifier operates on segmented cycles, errors in automatic cycle
        detection propagate into the final decision, coupling the two stages.
\end{itemize}

These challenges motivate a design that separates the difficulty of
\emph{classification} from the difficulty of automatic \emph{cycle detection},
so that the contribution of each can be measured rather than conflated.

% ---------------------------------------------------------------------
\section{Research Questions}
\label{sec:intro-rqs}

The thesis is organised around seven research questions:

\begin{description}
  \item[RQ1.] Can pose-based classifiers distinguish correct from false running
        technique when running cycles are manually or otherwise reliably
        segmented?
  \item[RQ2.] How much classification performance is lost when handmade cycles
        are replaced by automatic autocorrelation-based cycle detection?
  \item[RQ3.] Which motion signal is most reliable for automatic running-cycle
        detection?
  \item[RQ4.] Which model architecture is most suitable for classifying
        eight-phase pose sequences?
  \item[RQ5.] How do the feature representation and the number of frames sampled
        per cycle affect classification performance?
  \item[RQ6.] How does an image-based representation compare with a
        landmark-based representation under the same automatic
        (autocorrelation) pipeline?
  \item[RQ7.] How do cycle-level aggregation, decision-threshold selection, and
        pose/cycle quality affect the interpretation of the final results?
\end{description}

RQ1 establishes whether the task is solvable in principle; RQ2 quantifies the
cost of automation; RQ3--RQ5 examine the design choices internal to the
automatic pipeline; RQ6 addresses the representation question; and RQ7 concerns
how results should be measured and interpreted.

% ---------------------------------------------------------------------
\section{Contributions}
\label{sec:intro-contributions}

In addressing these questions, the thesis makes the following contributions:

\begin{enumerate}
  \item A complete markerless, pose-based pipeline for classifying running
        technique from in-the-wild internet videos, spanning pose estimation,
        landmark filtering and normalisation, motion-signal extraction,
        running-cycle detection, phase sampling, representation, deep-learning
        classification, and evaluation.
  \item A matched experimental design comprising a \emph{handmade upper bound},
        which characterises classification given reliable segmentation, and an
        \emph{autocorrelation-based automatic baseline}, which uses only
        information available from the video itself; the pairing isolates the
        effect of automatic cycle detection.
  \item An \emph{eight-phase} running-cycle representation that normalises for
        cadence so that cycles of differing duration become directly comparable.
  \item A set of controlled ablation studies on model architecture, detection
        signal, feature representation, and frames per cycle.
  \item Extended analyses comparing landmark-based and image-based
        representations on an identical cycle source, comparing cycle-level and
        video-level evaluation, examining decision-threshold selection, and
        relating pose and cycle quality to classification accuracy.
  \item A disciplined separation between \emph{historical development} results
        --- retained as method-evolution evidence --- and \emph{controlled
        benchmark} results obtained under a single fixed, leakage-free protocol,
        ensuring that development-stage milestones are not mistaken for
        benchmarks.
\end{enumerate}

% ---------------------------------------------------------------------
\section{Thesis Structure}
\label{sec:intro-structure}

The remainder of the thesis is organised as follows.
Chapter~\ref{ch:literature-review} reviews the relevant literature on running
technique, motion-analysis modalities, markerless pose estimation, running-cycle
detection, pose- and image-based classification, and evaluation protocols, and
identifies the research gap.
Chapter~\ref{ch:methodology} describes the proposed framework in detail,
including the dataset and leakage-free splitting protocol, pose estimation,
landmark filtering and normalisation, motion-signal extraction, the three
cycle-detection methods, eight-phase sampling, the landmark and image
representations, the classification models, the cycle- and video-level
prediction scheme, and threshold selection.
Chapter~\ref{ch:experimental-setup} specifies the experimental setup, including
training configurations, evaluation metrics, and the reporting protocol that
separates controlled from historical results.
Chapter~\ref{ch:results} reports the empirical results: dataset statistics, the
handmade upper bound, the automatic baseline, the ablation studies, the
historical version history, and the extended analyses.
Chapter~\ref{ch:discussion} interprets these results, locating where performance
is gained and lost across the pipeline.
Chapter~\ref{ch:conclusion} summarises the findings, restates the answers to the
research questions, and outlines directions for future work. Throughout, the
scope of the conclusions is kept commensurate with the dataset and protocol
used, and no claim of real-world deployment readiness is made.

% End of Chapter 1